\documentclass[10pt,a4paper]{amsart}
\usepackage[margin=19mm]{geometry}
\usepackage{amsmath,amssymb,mathtools}
\usepackage[scr=rsfs,cal=euler]{mathalfa}
\usepackage{xfrac}
\usepackage[unicode,hidelinks,bookmarks=false]{hyperref}
\newcommand{\ie}{i.e.\ }
\newcommand{\cf}{cf.\ }
\newcommand{\resp}{respectively\ }
\newcommand{\Subsection}{\subsection}
\providecommand{\nobreakdash}{\nobreak\hbox{-}}
\setlength{\emergencystretch}{2em}
\multlinegap0pt
\usepackage{mathtools}
\usepackage[all]{xy}
\usepackage{amssymb}
\usepackage{bbm}
\usepackage{caption}
\usepackage{enumerate}
\usepackage{xcolor}
\usepackage[shortlabels]{enumitem}
\newtheorem{lemma}{Lemma}
\newcommand{\A}{{\mathbb A}}
\newcommand{\Spec}{{\rm Spec \,}}
\newcommand{\sO}{{\mathcal O}}
\title{$\mathbb{A}^1$-fibration in algebraic geometry and $\mathbb{A}^1$-homotopy type.}
\author{Utsav Choudhury, Aritra Mandal, Biman Roy}
\date{Source: arXiv:2608.10966v1 (CC BY 4.0)}
\begin{document}
\maketitle

\noindent\emph{Source and licence.} $\mathbb{A}^1$-fibration in algebraic geometry and $\mathbb{A}^1$-homotopy type.. Version arXiv:2608.10966v1 (CC BY 4.0), distributed by arXiv under the Creative Commons Attribution 4.0 International licence, http://creativecommons.org/licenses/by/4.0/, as recorded in the arXiv licence metadata for this exact version and shown on the abstract page https://arxiv.org/abs/2608.10966v1. Full licence notice: \texttt{licenses/arxiv-2608.10966-CC-BY-4.0.txt}.\par\smallskip

\noindent\textit{Editorial scope.} Counted unit: the article's lemma lemma (source0.tex lines 1058--1063). The article's complete original proof of it is delivered with it (source0.tex lines 1064--1167, one \texttt{proof} environment of 104 lines). No mathematics is written, repaired or re-derived: the statement and the proof are the article's own lines, in the article's own notation, and no retained line is substituted or re-flowed.  No further source-local context is needed: every cross-reference of every retained line resolves inside the two delivered spans.  Nothing was omitted from any retained span, so the package carries no omission record.  No retained line contains a citation, so the wrapper carries no bibliography and no bibliography entry.  No retained line contains a figure environment, an image-inclusion command or drawing code, so no image is delivered and the package declares no asset dependency.  Editorial formatting only, in addition: an AMS \texttt{amsart} preamble that restates the article's own declarations and macros used by the retained lines, namely the environments \texttt{lemma} on separate counters, together with the standard packages those lines need.  The retained environments print in this document as Lemma 1 in order of appearance, whereas the article prints the same items with the numbering of its own full document.  The complete native source of the article as distributed in the arXiv e-print for this exact version is bundled unchanged as \texttt{sources/207371/source0.tex}.
\begin{lemma} \label{key lemma}
    Let $i:V \to U$ be a morphism between smooth varieties over $k$, which is either an open immersion or a closed immersion of schemes. Then the canonical morphism
    $$i_*: Sing_*(U)/Sing_*(V) \to Sing_*(U/V)$$
    is an isomorphism of simplicial sheaves.
    %Nisnevich local weak equivalence.
\end{lemma}

\begin{proof}
%    $Sing_*$ is a homotopy colimit of a simplicial diagram.
There is a commutative diagram of simplicial sheaves
  \[
\xymatrix{
Sing_*(V) \ar[r] \ar[d] & Sing_*(U) \ar[d] \\
\text{*} \ar[r] & Sing_*(U/V)
}
\]
This induces a morphism of simplicial sheaves
$$i_*: Sing_*(U)/Sing_*(V) \to Sing_*(U/V).$$
We will show that for every essentially smooth Henselian local scheme $\mathcal{O}$ over $k$, the morphism
$$i_*: (Sing_*(U)/Sing_*(V))(\mathcal{O}) \to Sing_*(U/V)(\mathcal{O})$$
is an isomorphism of simplicial sets, that is for every $n \geq 0$ the morphism between $n$-simplices
$$i_*: (Sing_*(U)/Sing_*(V))(\mathcal{O})_n \to Sing_*(U/V)(\mathcal{O})_n$$
is a bijection.
%n isomorphism.

Since over the Henselian local scheme, the presheaf section and sheaf section coincide, so
\begin{align*}
(Sing_*(U)/Sing_*(V))(\mathcal{O})_n & \cong Sing_*(U)(\sO)_n/Sing_*(V)(\sO)_n
%& \cong (Hom_{Sch/k}(\A^n_\sO, U)\coprod \{\bullet\})/\sim
\end{align*}
For any scheme $T$, the section $Sing_*(U)(T)_n/Sing_*(V)(T)_n$ is the pushout in the category of sets
 \[
\xymatrix{
Hom_{Sch/k}(\A^n_T, V) \ar[r] \ar[d] & Hom_{Sch/k}(\A^n_T, U) \ar[d] \\
\text{*} \ar[r] & Sing_*(U)(T)_n/Sing_*(V)(T)_n
}
\]
where $Sch/k$ is the category of $k$-schemes. So, $Sing_*(U)(T)_n/Sing_*(V)(T)_n$ is a pointed set canonically, denote the base point by $\bullet$. Also for the simplicial presheaf 
$$T \mapsto Sing_*(U)(T)/Sing_*(V)(T),$$
the restriction map
$$Sing_*(U)(T)/Sing_*(V)(T) to Sing_*(U)(T^\prime)/Sing_*(V)(T^\prime)$$
induced by a morphism $T^\prime \to T$ of $k$-schemes, respect the base points.

%$(Sing_*(U)/Sing_*(V))_n(\mathcal{O})$ is a pointed set canonically, denote it by $\bullet$. 
We will prove that the morphism  $i_*$ is a bijection between the pointed sets
$$i_*: ((Sing_*(U)/Sing_*(V))(\mathcal{O})_n, \bullet) \to (Sing_*(U/V)(\mathcal{O})_n, i_*(\bullet)),$$
for every Henselian local scheme $\sO$.

%$$

%Fix a $k$-point $v \in V(k)$. The point $v$ determines an element in the simplicial set $(Sing_*(U)/Sing_*(V))(\mathcal{O})$ (which we will denote by $v$ again),
%and $Sing_*(U/V)(\mathcal{O})$ 
%by taking the composition of the morphisms
%$$\mathbb{A}^n_{\mathcal{O}} \to \Spec\ k \xrightarrow{v} V.$$
%The elements $v \in (Sing_*(U)/Sing_*(V))(\mathcal{O})$ and $i_*(v) \in Sing_*(U/V)(\mathcal{O})$ make both the simplicial sets to be pointed. 
%For each $n$, fix a distinguished element $c: \mathbb{A}^n_{\mathcal{O}} \to V$ given by the constant map to a $k$ -point of $V$; that makes both the simplicial sets to the pointed.

\textbf{Injectivity:} Suppose, $\phi, \psi \in (Sing_*(U)/Sing_*(V))(\mathcal{O})_n$ and 
$$i_*(\phi) = i_*(\psi) \in Sing_*(U/V)(\mathcal{O}).$$ 
For presheaf and the associated sheaf, the sections over $\mathcal{O}$ are same; thus $\phi, \psi \in (Sing_*(U)/Sing_*(V))^{\text{pre}}(\mathcal{O})$. So the sections are given by the morphisms $\phi, \psi: \mathbb{A}^n_{\mathcal{O}} \to U$ (again, denote by $\phi$ and $\psi$). 
%We denote the images of $\phi$ and $\psi$ in the quotient $U(\mathbb{A}^n_{\mathcal{O}})/V(\mathbb{A}^n_{\mathcal{O}})$ by $[\phi]$ and $[\psi]$ respectively.
% such that their images $[\phi], [\psi] \in U(\mathbb{A}^n_{\mathcal{O}})/V(\mathbb{A}^n_{\mathcal{O}})$ are same. 
Since $i_*(\phi) = i_*(\psi) \in (U/V)(\mathbb{A}^n_{\mathcal{O}})$, so there is a Nisnevich covering $f: W \to \mathbb{A}^n_{\mathcal{O}}$ such that 
$$\phi \circ f = \psi \circ f \in (U/V)^{\text{pre}}(W) = U(W)/V(W).$$
This means either $\phi \circ f, \psi \circ f: W \to V \hookrightarrow U$ or $\phi \circ f = \psi \circ f: W \to U$.
Since $Hom(-, U)$ is a Nisnevich sheaf, so for the second case we conclude $\phi = \psi: \mathbb{A}^n_{\mathcal{O}} \to U$. For the first case, since $f$ is a covering, in particular surjective; so the images of $\phi$ and $\psi$ are in $V$, thus both the morphisms $\phi$ and $\psi$ factor as
$$\phi: \mathbb{A}^n_{\mathcal{O}} \to V \hookrightarrow U \text{ and } \psi: \mathbb{A}^n_{\mathcal{O}} \to V \hookrightarrow U.$$
Hence 
$$\phi = \psi \in (Sing_*(U)/Sing_*(V))^{\text{pre}}(\mathcal{O}) = U(\mathbb{A}^n_{\mathcal{O}})/V(\mathbb{A}^n_{\mathcal{O}}).$$
So $i_*$ is injective.

\textbf{Surjectivity: } Suppose, $\alpha \in Sing_*(U/V)(\mathcal{O})_n$. So $\alpha \in (U/V)(\mathbb{A}^n_{\mathcal{O}})$. Thus there is a Nisnevich covering $W \to \mathbb{A}^n_{\mathcal{O}}$ such that $\alpha$ is given by an element $\phi \in (U/V)^{\text{pre}}(W)$ (that is, $\phi: W \to U$ is a morphism) such that there is a Nisnevich covering $\Tilde{W} \to W \times_{\mathbb{A}^n_{\mathcal{O}}} W$ such that 
$$\text{pr}_1^*(\phi)|_{\Tilde{W}} = \text{pr}_2^*(\phi)|_{\Tilde{W}} \in (U/V)^{\text{pre}}(\Tilde{W}) = U(\Tilde{W})/V(\Tilde{W}).$$
This means either 
$$\text{pr}_1^*(\phi)|_{\Tilde{W}}, \text{pr}_2^*(\phi)|_{\Tilde{W}} : \Tilde{W} \to V \hookrightarrow U$$
or
$$\text{pr}_1^*(\phi)|_{\Tilde{W}} = \text{pr}_2^*(\phi)|_{\Tilde{W}} : \Tilde{W} \to U.$$
For the second case, since $Hom(-, U)$ is a Nisnevich sheaf and the morphism $\Tilde{W} \to W \times_{\mathbb{A}^n_{\mathcal{O}}} W$ is a Nisnevich covering; so 
$$\text{pr}_1^*(\phi) = \text{pr}_2^*(\phi) : W \times_{\mathbb{A}^n_{\mathcal{O}}} W \to U$$
and again as $W \to \mathbb{A}^n_{\mathcal{O}}$ is a Nisnevich covering; so $\phi$ can be lifted to a morphism $\bar{\phi}: \mathbb{A}^n_{\mathcal{O}} \to U$, which proves surjectivity in this case. \par
For the first case, if $\text{pr}_1^*(\phi)|_{\Tilde{W}}, \text{pr}_2^*(\phi)|_{\Tilde{W}} : \Tilde{W} \to V$, then we claim that the morphism $\phi$ factors through $V$. Indeed, the composition of morphisms
$$\Tilde{W} \to W \times_{\mathbb{A}^n_{\mathcal{O}}} W \xrightarrow{\text{pr}_1} W$$
and 
$$\Tilde{W} \to W \times_{\mathbb{A}^n_{\mathcal{O}}} W \xrightarrow{\text{pr}_2} W$$
are Nisnevich coverings of $W$ and the images of $\text{pr}_1^*(\phi)|_{\Tilde{W}}$ and $\text{pr}_2^*(\phi)|_{\Tilde{W}}$ are in $V$, so the image of $\phi$ is also in $V$. Now we claim that $\alpha = i_*(\bullet) \in (U/V)(\mathbb{A}^n_{\mathcal{O}})$, where $\bullet$ is the distinguished element in $(Sing_*(U)/Sing_*(V))(\mathcal{O})$ mentioned before. This is equivalent to show that 
$$\alpha|_W = i_*(\bullet)|_W \in (U/V)(W),$$
where $W \to \mathbb{A}^n_{\mathcal{O}}$ is the Nisnevich covering and $-|_W$ denotes the restriction maps, induced by $W \to \A^n_\sO$, for both the presheaf $(U/V)^{\text{pre}}$, given by 
$$(U/V)^{\text{pre}}(T) = U(T)/V(T)$$
and the associated Nisnevich sheaf $U/V$.
% the restriction map of the quotient sheaf $U/V$
%$$-|_W:(U/V)(\mathbb{A}^n_{\mathcal{O}}) \to (U/V)(W).$$
%The quotient sheaf $U/V$ is the Nisnevich sheaf associated to the presheaf $(U/V)^{\text{pre}}$, which is given by 
%$$(U/V)^{\text{pre}}(T) = U(T)/V(T).$$
Now, $\alpha|_W = \eta(\phi)$. Consider the commutative diagram
 \[
\xymatrix{
U(\A^n_\sO)/V(\A^n_\sO) \ar[r]^\eta \ar[d]^{|_W} & (U/V)(\A^n_\sO) \ar[d]^{|_W} \\
U(W)/V(W) \ar[r]_\eta & (U/V)(W)
}
\]
%$i_*(\bullet)|_W = \eta(\bullet|_W)$, 
where 
$$\eta: (U/V)^{\text{pre}} \to U/V$$ 
is the sheafification morphism and in the square, the left vertical morphism respects the base point. So,
$$i_*(\bullet)|_W = \eta(\bullet|_W).$$
Since the image of $\phi$ is contained in $V$, so $\phi = \bullet|_W \in U(W)/V(W)$, so 
$$\alpha|_W =\eta(\bullet|_W)= i_*(\bullet)|_W \in (U/V)(W)$$ 
and consequently, $\alpha = i_*(\bullet)$. This proves the surjectivity. 

This completes the proof.
\end{proof}

\end{document}
